\documentclass[11pt]{article}
\usepackage[margin=1in]{geometry}
\usepackage{amsmath,amssymb,amsthm}
\usepackage{graphicx}
\usepackage[hidelinks]{hyperref}

\newtheorem{theorem}{Theorem}
\newtheorem{lemma}{Lemma}
\newtheorem{remark}{Remark}

\newcommand{\DPP}{Dunford--Pettis property}
\newcommand{\B}{\mathcal B}
\newcommand{\PiOne}{\Pi_1}
\newcommand{\piOne}{\pi_1}

\title{Hadamard Block Counterexamples to the \(1\)-DPP of \(\ell_p\)}
\author{FA Banach run \texttt{fa\_banach\_001}, agent \texttt{agent\_lane\_05}}
\date{2026-06-16}

\begin{document}
\maketitle

\section*{Status}

This packet records a candidate counterexample, likely valid. It gives a
negative answer to a question left open by Karn and Sinha in arXiv:1207.3620:
whether \(\ell_p\), \(1<p<\infty\), has the \(1\)-Dunford--Pettis property.

\section*{Source Passage}

After proving that \(\ell_p\) does not have the \(r\)-DPP for any \(r>1\), and
that it does not have several mixed \((r,s)\)-DPP properties with \(r>1\) or
\(s>1\), Karn and Sinha leave the endpoint \(r=s=1\) unsettled.

\begin{center}
\includegraphics[width=0.98\linewidth]{figures/open_problem_crop.png}
\end{center}

\section*{Transcription of the Question}

The note on page 11 states:
\[
\text{``We have not been able to settle whether for }1<p<\infty,
\ \ell_p\text{ has the }1\text{-DPP.''}
\]

\section*{Definitions Used}

We use the characterization from Karn and Sinha's Theorem 3.7
\cite{KarnSinha}: a Banach space \(X\) has the \(1\)-Dunford--Pettis property
if and only if
\[
  \Pi_1(X,\ell_1)=\B(X,\ell_1).
\]
Thus, to disprove the \(1\)-DPP for \(X=\ell_p\), it is enough to construct a
bounded operator \(T:\ell_p\to\ell_1\) which is not absolutely \(1\)-summing.

For a bounded operator \(S:X\to Y\), the absolutely \(1\)-summing norm is
denoted by \(\piOne(S)\). We use the standard finite-sequence lower test
\[
  \piOne(S)\ge
  \frac{\sum_i \|Sx_i\|}
       {\sup_{\varphi\in B_{X^*}}\sum_i |\varphi(x_i)|}.
\]

\section*{The Counterexample}

\begin{theorem}
For every \(1<p<\infty\), the space \(\ell_p\) does not have the
\(1\)-Dunford--Pettis property.
\end{theorem}

\section*{Proof Intuition}

The obstruction is finite-dimensional.  A Hadamard matrix \(H_N\), viewed as
an operator \(\ell_p^N\to\ell_1^N\), has small operator norm because
\(\|H_Nx\|_1\le N\|x\|_2\).  But its absolutely \(1\)-summing norm is large:
on the \(N\) unit vectors, all columns have \(\ell_1\)-norm \(N\), while the
weak \(\ell_1\)-size of the unit vectors in \(\ell_p^N\) is only \(N^{1/p}\).
Therefore \(\piOne(H_N)/\|H_N\|\to\infty\).

The infinite counterexample is a block diagonal sum of these finite matrices.
The block scalars make the operator norms summable in the correct
\(\ell_{p'}\)-sense, so the resulting operator \(\ell_p\to\ell_1\) is bounded.
At the same time, the \(1\)-summing norm on the individual blocks still tends
to infinity.  A \(1\)-summing operator cannot have restrictions with
unbounded \(1\)-summing norm, so the block sum is bounded but not
absolutely \(1\)-summing.

\begin{proof}
Let \(p'\) be the conjugate exponent of \(p\).  For each \(N=2^j\), let
\(H_N\) be the Sylvester Hadamard matrix, so its entries are \(\pm 1\) and
\[
  H_N^tH_N=NI_N.
\]
Regard \(H_N\) as an operator \(\ell_p^N\to\ell_1^N\).

First we estimate its operator norm.  For every \(x\in\mathbb R^N\) or
\(\mathbb C^N\),
\[
  \|H_Nx\|_1
  \le N^{1/2}\|H_Nx\|_2
  = N\|x\|_2.
\]
Therefore
\[
  \|H_N:\ell_p^N\to\ell_1^N\|
  \le B_N :=
  \begin{cases}
    N, & 1<p\le 2,\\
    N^{3/2-1/p}, & 2\le p<\infty,
  \end{cases}
\]
using \(\|x\|_2\le \|x\|_p\) for \(p\le2\) and
\(\|x\|_2\le N^{1/2-1/p}\|x\|_p\) for \(p\ge2\).

Next we lower-bound \(\piOne(H_N)\).  Test on the unit vectors
\(e_1,\ldots,e_N\) of \(\ell_p^N\).  Each column \(H_Ne_i\) has
\(\ell_1\)-norm \(N\), and
\[
  \sup_{\|a\|_{p'}\le1}\sum_{i=1}^N |\langle a,e_i\rangle|
  = \sup_{\|a\|_{p'}\le1}\|a\|_1
  = N^{1/p}.
\]
Hence
\[
  \piOne(H_N)\ge \frac{N^2}{N^{1/p}}=N^{2-1/p}.
\]
It follows that
\[
  \frac{\piOne(H_N)}{B_N}\ge R_N :=
  \begin{cases}
    N^{1-1/p}=N^{1/p'}, & 1<p\le2,\\
    N^{1/2}, & 2\le p<\infty.
  \end{cases}
\]
In either case \(R_N\to\infty\) along \(N=2^j\).

Now take \(N_j=2^j\), \(R_j=R_{N_j}\), and \(B_j=B_{N_j}\).  Set
\[
  c_j=R_j^{-1/2},\qquad A_j=\frac{c_j}{B_j}H_{N_j}:
  \ell_p^{N_j}\to\ell_1^{N_j}.
\]
Then \(\|A_j\|\le c_j\).  Since \(R_j\) grows exponentially, the scalar
sequence \((c_j)\) belongs to \(\ell_{p'}\).  Identify
\[
  \ell_p \cong \left(\bigoplus_{j=1}^\infty \ell_p^{N_j}\right)_{\ell_p},
  \qquad
  \ell_1 \cong \left(\bigoplus_{j=1}^\infty \ell_1^{N_j}\right)_{\ell_1}.
\]
Define the block diagonal operator
\[
  T((x_j)_{j=1}^\infty)=(A_jx_j)_{j=1}^\infty.
\]
For \(x=(x_j)\in\ell_p\), Holder's inequality gives
\[
  \|Tx\|_1
  =\sum_j \|A_jx_j\|_1
  \le \sum_j c_j\|x_j\|_p
  \le \|(c_j)\|_{p'} \left(\sum_j \|x_j\|_p^p\right)^{1/p}.
\]
Thus \(T:\ell_p\to\ell_1\) is bounded.

However,
\[
  \piOne(A_j)
  =\frac{c_j}{B_j}\piOne(H_{N_j})
  \ge c_jR_j
  =R_j^{1/2}\longrightarrow\infty.
\]
If \(T\) were absolutely \(1\)-summing, then for the coordinate inclusion
\(J_j:\ell_p^{N_j}\to\ell_p\) and coordinate projection
\(P_j:\ell_1\to\ell_1^{N_j}\), the ideal property would imply
\[
  \piOne(A_j)=\piOne(P_jTJ_j)\le \piOne(T)
\]
for every \(j\), a contradiction.  Therefore \(T\) is bounded but not
absolutely \(1\)-summing.

By Karn and Sinha's characterization
\(\Pi_1(\ell_p,\ell_1)=\B(\ell_p,\ell_1)\) if and only if \(\ell_p\) has the
\(1\)-DPP.  The operator \(T\) shows this equality fails, and hence
\(\ell_p\) does not have the \(1\)-DPP.
\end{proof}

\section*{Verification Notes}

The proof is analytic.  The included script
\texttt{code/verify\_hadamard\_bounds.py} verifies the finite Hadamard
identities \(H_N^tH_N=NI_N\) and the column \(\ell_1\)-norms for small
Sylvester matrices, then prints the lower and upper estimates used above for
sample values of \(p\).  It is a sanity check, not a substitute for the proof.

\section*{Novelty and Literature Check}

The run indexes were searched for arXiv:1207.3620, the paper title,
``operator summability'', ``operator p-summability'', ``1-DPP'',
``Pi\_1(ell\_p,ell\_1)'', and close variants.  No prior packet or attempt in
the run addressed this endpoint question.

The local parsed arXiv corpus was searched for later citations to Karn--Sinha
and for exact phrases around \(1\)-DPP, \(\ell_p\), and absolutely
\(1\)-summing operators into \(\ell_1\).  Several related papers cite the
source, but this bounded pass found no explicit later answer to the note.
External web searches for exact phrases around ``ell\_p 1-DPP'', ``l\_p
1-DPP An operator summability'', ``Pi\_1 ell\_p ell\_1 1-summing'', and
``absolutely 1-summing ell\_p ell\_1'' likewise surfaced no explicit later
resolution.  Novelty confidence is moderate: the result appears new to the
run memory and to this bounded search, but the literature check is not
exhaustive.

\section*{Limitations}

This answers the endpoint question for the classical spaces \(\ell_p\),
\(1<p<\infty\).  It does not attempt to classify all spaces for which every
operator into \(\ell_1\) is absolutely \(1\)-summing, nor does it address the
separate Section 4 question in the source paper about \(\ell_p^t\)-completeness
of the operator ideal \(Lt_p^\cdot(X,Y)\).

\section*{Human Review Recommendation}

Reviewers should check three points first: the finite-dimensional
\(\piOne(H_N)\) lower bound obtained from the unit vectors, the two-case
operator-norm estimate through \(\ell_2\), and the weighted block-diagonal
argument showing boundedness while the block restrictions have
\(\piOne\)-norms tending to infinity.

\begin{thebibliography}{9}

\bibitem{KarnSinha}
A. K. Karn and D. P. Sinha,
\emph{An operator summability of sequences in Banach spaces},
Glasgow Math. J. \textbf{56} (2014), 427--437; arXiv:1207.3620.

\bibitem{DJT}
J. Diestel, H. Jarchow, and A. Tonge,
\emph{Absolutely Summing Operators},
Cambridge Studies in Advanced Mathematics 43, Cambridge University Press,
1995.

\end{thebibliography}

\end{document}
